\subsubsection{分布}\label{subsec:distributions-definition}

现在定义分布空间. 如引言所述, 分布的定义构造得非常巧妙, 以实现两个看似矛盾的目标. 希望有一个"函数"的广义概念, 其中包含高度奇异或"粗糙"的对象. 同时, 又希望能够定义这些对象的任意阶"导数".

\begin{Definition}
\label{def:distribution}
\textbf{分布}或\textbf{广义函数}是从 $\mathcal{D}(\Omega)$ 到 $\mathbb{R}$ 的线性映射 $\phi\mapsto(f,\phi)$, 它在以下意义下连续: 若 $\phi_n\to\phi$ 在 $\mathcal{D}(\Omega)$ 中成立, 则 $(f,\phi_n)\to(f,\phi)$. 所有分布构成的集合称为 $\mathcal{D}^{\prime}(\Omega)$.
\end{Definition}

\begin{Example}
\label{ex:regular-distribution}
对 $\Omega$ 上任意连续函数 $f$, 令
\begin{equation}\label{eq:regular-distribution-pairing}
(f,\phi)=\int_\Omega f(\bm{x})\phi(\bm{x})\,\mathrm{d}\bm{x}.
\end{equation}
即可将其与一个广义函数等同. 映射的连续性由关于一致收敛函数序列的积分极限的熟知定理得出. 事实上, 即使 $f$ 仅局部可积, Lebesgue 控制收敛定理仍允许作出相同断言.
\end{Example}

\begin{Example}
\label{ex:dirac-delta}
当然, 有许多广义函数并不对应于通常意义下的"函数". 最重要的例子称为 Dirac delta 函数. 假设 $\Omega$ 包含原点, 定义
\begin{equation}\label{eq:dirac-delta-pairing}
(\delta,\phi)=\phi(0).
\end{equation}
泛函的连续性来自试验函数序列的收敛蕴含逐点收敛这一事实.

容易证明, 不存在满足 \eqref{eq:dirac-delta-pairing} 的连续函数 (参见 Problem~\ref{ex:delta-not-continuous-function}).
\end{Example}

\begin{Remark}
\label{rem:generalized-function-point-values}
像 delta 函数这样的广义函数, 并不像普通函数那样取"值". 尽管如此, 习惯上仍使用普通函数的语言, 谈论"广义函数 $\delta(\bm{x})$",\footnote{对于那些认为这种说法令人厌恶的朋友, 我们表示歉意 --- 即使对普通函数, 他们也认为这种说法令人厌恶!} 即使代入某个具体的 $\bm{x}$ 并没有意义. 也将用 $\int_\Omega\delta(\bm{x})\phi(\bm{x})\,\mathrm{d}\bm{x}$ 表示 $(\delta,\phi)$.
\end{Remark}

\begin{Example}
\label{ex:point-derivative-functional}
对任意多重指标 $\alpha$, 映射 \[\phi\mapsto D^\alpha\phi(0)\] 是广义函数.
\end{Example}

\begin{Example}
\label{ex:surface-charge-distribution}
其他奇异分布包括面电荷这类来自物理学的例子. 若 $S$ 是 $\mathbb{R}^3$ 中光滑的二维曲面, 且 $q:S\to\mathbb{R}$ 可积, 则对 $\phi\in\mathcal{D}(\mathbb{R}^3)$ 定义\[(q,\phi)=\int_Sq(\bm{x})\phi(\bm{x})\,\mathrm{d}a(\bm{x})\]其中 $\mathrm{d}a(\bm{x})$ 表示关于 $S$ 上的曲面面积积分.
\end{Example}

\begin{Example}
\label{ex:curve-current-distribution}
沿曲线 $\mathcal{C}\subset\mathbb{R}^3$ 流动的电流是向量值分布的一个例子. 若 $\bm{j}:\mathcal{C}\to\mathbb{R}^3$ 可积, 则对 $\bm{\phi}\in\mathcal{D}(\mathbb{R}^3)^3$ 定义\[(\bm{j},\bm{\phi})=\int_{\mathcal{C}}\bm{j}(\bm{x})\cdot\bm{\phi}(\bm{x})\,\mathrm{d}\sigma(\bm{x})\]其中 $\mathrm{d}\sigma(\bm{x})$ 表示关于 $\mathcal{C}$ 上的弧长积分.
\end{Example}

\begin{Remark}
\label{rem:complex-distribution-convention}
当然, 可以用与实值分布相同的方式定义复值分布; 不过在这种情况下, 习惯采用约定
\begin{equation}\label{eq:complex-distribution-conjugate-pairing}
(f,\phi)=\int_\Omega\overline{f(\bm{x})}\phi(\bm{x})\,\mathrm{d}\bm{x}
\end{equation}
代替 \eqref{eq:regular-distribution-pairing}; 这样, 广义函数与试验函数的配对就具有与 Hilbert 空间 $L^2(\Omega)$ 中的内积相同的形式.\footnote{Hilbert 空间理论中最古老的问题之一, 是内积中的复共轭应放在第一个因子上还是第二个因子上. 这里的约定为物理学家广泛采用. 纯数学家倾向于采用相反的约定.}
\end{Remark}

分布的一个重要性质是, 它们局部具有"有限阶".

\begin{Lemma}
\label{lem:distribution-local-finite-order}
设 $f\in\mathcal{D}^{\prime}(\Omega)$, 且 $K$ 为 $\Omega$ 的紧子集. 那么存在 $n\in\mathbb{N}$ 和常数 $C$, 使对每个支集包含于 $K$ 的 $\phi\in\mathcal{D}(\Omega)$, 都有
\begin{equation}\label{eq:distribution-local-finite-order-bound}
|(f,\phi)|\leq C\sum_{|\alpha|\leq n}\max_{\bm{x}\in K}|D^\alpha\phi(\bm{x})|
\end{equation}

\end{Lemma}

\begin{Proof}
反设不成立. 那么对每个 $n$, 都存在 $\psi_n$, 使\[|(f,\psi_n)|>n\sum_{|\alpha|\leq n}\max_{\bm{x}\in K}|D^\alpha\psi_n(\bm{x})|.\]令 $\phi_n:=\psi_n/|(f,\psi_n)|$. 那么在 $\mathcal{D}(\Omega)$ 中 $\phi_n\to0$, 但 $(f,\phi_n)\equiv1$. 这产生矛盾, 证明完成.
\end{Proof}

以一些直接的定义结束本小节.

\begin{Definition}
\label{def:distribution-linear-operations}
对分布 $f$ 和 $g$ 及实数 $\alpha\in\mathbb{R}$, 令
\begin{equation}\label{eq:distribution-addition}
(f+g,\phi)=(f,\phi)+(g,\phi),
\end{equation}

\begin{equation}\label{eq:distribution-scalar-multiplication}
(\alpha f,\phi)=(f,\alpha\phi).
\end{equation}
(若允许 $\alpha$ 为复数, 则将 \eqref{eq:distribution-scalar-multiplication} 的右端改为 $(f,\overline\alpha\phi)$.)
\end{Definition}

\begin{Remark}
\label{rem:distribution-products}
一般不能定义两个广义函数的乘积 (参见 Problems~\ref{ex:distribution-multiplication-not-continuous}, \ref{ex:delta-approximation-square-divergence}). 然而, 可以定义分布与光滑函数的乘积.
\end{Remark}

\begin{Definition}
\label{def:distribution-smooth-product}
对任意函数 $a\in C^\infty(\Omega)$, 定义
\begin{equation}\label{eq:distribution-smooth-function-product}
(af,\phi)=(f,a\phi).
\end{equation}

\end{Definition}

若将函数 $f(x)$ 的图像平移 $h$, 则得到函数 $f(x-h)$ 的图像, 即 $x$ 平移 $-h$. 这一点可推广到 $\mathbb{R}^m$ 上的分布.

\begin{Definition}
\label{def:distribution-affine-change}
设 $\bm{U}(\bm{x})=\bm{A}\bm{x}+\bm{b}$ 为 $\mathbb{R}^m$ 中的非奇异线性变换, 设 $\bm{U}^{-1}(\bm{y})=\bm{A}^{-1}(\bm{y}-\bm{b})$ 为逆变换. 那么令
\begin{equation}\label{eq:distribution-affine-change-variables}
(\bm{U}f,\phi)=|\det A|\,(f(\bm{x}),\phi(\bm{U}(\bm{x}))).
\end{equation}

\end{Definition}

这一定义的动机来自以下形式计算:\[
\begin{aligned}
(\bm{U}f,\phi)&=(f(\bm{U}^{-1}(\bm{x})),\phi(\bm{x}))\\
&=\int_{\mathbb{R}^m}f(\bm{U}^{-1}(\bm{x}))\phi(\bm{x})\,\mathrm{d}\bm{x}\\
&=|\det A|\int_{\mathbb{R}^m}f(\bm{y})\phi(\bm{U}(\bm{y}))\,\mathrm{d}\bm{y}.
\end{aligned}
\](这里作了代换 $\bm{x}=\bm{U}(\bm{y})$.)

\begin{Example}
\label{ex:translated-dirac-delta}
平移 $\delta(\bm{x}-\bm{x}_0)$ 定义为
\begin{equation}\label{eq:translated-dirac-delta-pairing}
(\delta(\bm{x}-\bm{x}_0),\phi(\bm{x}))=(\delta(\bm{x}),\phi(\bm{x}+\bm{x}_0))=\phi(\bm{x}_0).
\end{equation}

\end{Example}

\begin{Remark}
\label{rem:distribution-symmetry}
利用这一定义, 可以定义广义函数的对称性; 例如, 若 $f(-\bm{x})=f(\bm{x})$, 即 $(f(\bm{x}),\phi(\bm{x}))=(f(\bm{x}),\phi(-\bm{x}))$, 则 $f$ 为偶的.
\end{Remark}
